\begin{pastexam}[label={px:ntu-109-comm-dsa-35}]{臺大電信所 109 學年度（B）第 35 題}
（單選；每題 2 分，答錯不倒扣）

Consider the following two sequences:\\
3, 2, 1, 7, 9, 8, 6, 8, 0, 9\\
8, 3, 7, 8, 1, 9, 8, 3, 9, 0\\
The length of the longest common subsequence is:
\begin{enumerate}[label=(\Alph*)]
\item 3
\item 5
\item 4
\item 2
\end{enumerate}
\end{pastexam}

\begin{pastsoln}
答案是 (B)。 注意選項不是由小到大排，$5$ 是 (B)。

令 $X = \langle 3, 2, 1, 7, 9, 8, 6, 8, 0, 9\rangle$、$Y = \langle 8, 3, 7, 8, 1, 9, 8, 3, 9, 0\rangle$，
$c[i, j]$ 是 $X$ 前 $i$ 個和 $Y$ 前 $j$ 個的 LCS 長度（資結書第~18.3~節）：
\[
c[i, j] =
\begin{cases}
0, & i = 0 \text{ 或 } j = 0,\\
c[i-1, j-1] + 1, & x_i = y_j,\\
\max(c[i-1, j],\ c[i, j-1]), & x_i \ne y_j .
\end{cases}
\]
逐列填表（粗體是下面回溯時取到的相等格）：
\begin{center}
\small
\begin{tabular}{cc|ccccccccccc}
\toprule
 & $j$ & 0 & 1 & 2 & 3 & 4 & 5 & 6 & 7 & 8 & 9 & 10 \\
$i$ & $x_i \backslash y_j$ & & $8$ & $3$ & $7$ & $8$ & $1$ & $9$ & $8$ & $3$ & $9$ & $0$ \\ \midrule
$0$ & & 0 & 0 & 0 & 0 & 0 & 0 & 0 & 0 & 0 & 0 & 0 \\
$1$ & $3$ & 0 & 0 & \textbf{1} & 1 & 1 & 1 & 1 & 1 & 1 & 1 & 1 \\
$2$ & $2$ & 0 & 0 & 1 & 1 & 1 & 1 & 1 & 1 & 1 & 1 & 1 \\
$3$ & $1$ & 0 & 0 & 1 & 1 & 1 & \textbf{2} & 2 & 2 & 2 & 2 & 2 \\
$4$ & $7$ & 0 & 0 & 1 & 2 & 2 & 2 & 2 & 2 & 2 & 2 & 2 \\
$5$ & $9$ & 0 & 0 & 1 & 2 & 2 & 2 & \textbf{3} & 3 & 3 & 3 & 3 \\
$6$ & $8$ & 0 & 1 & 1 & 2 & 3 & 3 & 3 & \textbf{4} & 4 & 4 & 4 \\
$7$ & $6$ & 0 & 1 & 1 & 2 & 3 & 3 & 3 & 4 & 4 & 4 & 4 \\
$8$ & $8$ & 0 & 1 & 1 & 2 & 3 & 3 & 3 & 4 & 4 & 4 & 4 \\
$9$ & $0$ & 0 & 1 & 1 & 2 & 3 & 3 & 3 & 4 & 4 & 4 & \textbf{5} \\
$10$ & $9$ & 0 & 1 & 1 & 2 & 3 & 3 & 4 & 4 & 4 & 5 & 5 \\
\bottomrule
\end{tabular}
\end{center}
$c[10, 10] = 5$。 從右下角回溯：相等就往左上並記下這個字元，不相等就往比較大的那一格走
（一樣大時往上）。 取到的相等格是 $(9, 10)$、$(6, 7)$、$(5, 6)$、$(3, 5)$、$(1, 2)$，反過來就是
一個 LCS：$\langle 3, 1, 9, 8, 0\rangle$。 對照原序列：在 $X$ 是第 $1, 3, 5, 6, 9$ 個，在 $Y$ 是
第 $2, 5, 6, 7, 10$ 個，順序都對。
\end{pastsoln}
